For Proof of Reflection to work effectively, there must be some method
of comparing and converting the \emph{work} done by reflecting chains.
Earlier, we simply \emph{set} a ratio between Chain L blocks and Chain R
blocks, but that isn't suitable for a production network.

How can we design a system that allows for sensible comparisons between
Proofs of Work that use different hashing algorithms?

The core problem is that the work done on different chains is measured
using different units -- and those units aren't convertible. This
applies to blockchains that use the same hashing algorithm, too, since
there may be different costs and factors that are implicit in the mining
of each of those chains. One hash is not necessarily worth exactly one
hash in different contexts.

Revisiting \emph{qualitative conversion} from \autoref{sec:por-step4}:
the way we will analyze (and criticize) potential solutions is via the
concept of IGCs -- \{Idea, Goal, Context\} triples\footnote{IGCs are a
  method of structuring ideas (solutions to problems) so that they can
  be effectively analyzed and criticized. They're also an introduction
  to the thinking methods and techniques of \emph{Critical Fallibilism}.
  See \url{https://curi.us/2387-igcs}.}. The \emph{idea} component
contains \emph{the method} of conversion such that the the goal is
satisfied in the given context. The three elements of an IGC are all of
the components that are required to evaluate criticisms and thus
determine whether the IGC succeeds or fails.

In order to convert otherwise unconvertible units, one must define a
suitable goal and context for that conversion to make sense. For
example: \emph{is a cucumber longer than it is green?}\footnote{This
  example is from Eli Goldratt's \emph{The Choice} (2008).} That
question doesn't make sense because we can't convert between length and
color. However, consider the situation where you want to win a cucumber
competition, and points are awarded for a cucumber based on both its
consistency of color and its length (and nothing else). Based on the
specific rules, you could figure out a way to convert both color and
length into points. This would help you pick the best of your cucumbers
to enter into the competition, and your developing of that method means
that you also now have a way to convert color and length via whatever
relationship you came up with. Depending on the specific rules, you
could now say things like \emph{1cm of length is worth 3 blemishes}. In
order to make sense of \emph{converting between a cucumber's color and
length}, you need both the \emph{context} of the competition's rules and
also the \emph{goal} of maximizing the number of points\footnote{NB:
  Other \{goal, context\} pairs could work, too, and would have
  different methods of conversion.}.

In the case of \emph{Proof of Reflection}, we need to define a suitable
goal and context to enable this conversion, and come up with ideas for
how to do that conversion. Our goal is straightforward, and the same as
other consensus mechanisms: we want a blockchain system that is as
difficult as possible to attack. The context is the architectures of
both the PoR implementation and the blockchains in question, plus the
rest of the world (including attackers). More specifically: the
difficulty of an attack against a blockchain network is typically
analyzed from the context of an attacker's \emph{risk vs reward} where
the attacker has a goal of \emph{profit}.

Two ideas follow that create two distinct IGCs.
